% Part:sets-functions-relations
% Chapter: size-of-sets
% Section: enumerations-alt

\documentclass[../../../include/open-logic-section]{subfiles}

\begin{document}

\olfileid{sfr}{siz}{enm-alt}

\olsection{枚举与可数集}

\begin{editorial}
  本节按照集合论中的方式，将枚举定义为从 $\Nat$（或其初始段）到集合的双射。因此，它与 \olref[enm]{sec} 节中的定义略有冲突，并且重复了那里的所有例子。同时，本节也比那一节更简练。
\end{editorial}

我们可以通过直接枚举其元素来确定一个有限集合。我们在如下定义集合时就是这样做的：
\[
  A = \{a_1, a_2, \ldots, a_n\}.
\]
假设元素 $a_1$, \dots, $a_n$ 互不相同，这就给出了 $A$ 与前 $n$ 个自然数 $0$, \dots, $n-1$ 之间的一个双射。反过来，由于每个有限集都只有有限多个元素，所以每个有限集都能与自然数的一个初始段建立这样的对应关系。换句话说，如果 $A$ 是有限的，则在 $A$ 与 $\{0, \dots, n-1\}$ 之间存在一个双射，其中 $n$ 是~$A$ 的元素个数。

如果我们容许某些种类的无穷集，那么我们也会允许对一些无穷集进行枚举。精确地说，一个无穷集是通过它与整个~$\Nat$ 之间的双射来枚举的。

\begin{defn}[枚举（集合论）]
一个集合 $A$ 的\emph{枚举}，是指以自然数的某个初始段 $\{0,
1, \ldots, n\}$ {或}整个自然数集~$\Nat$ 为定义域、以 $A$ 为值域的双射。
\end{defn}

\begin{explain}
对\emph{枚举}一词的这种用法有其直观基础。因为说我们枚举了一个集合 $A$，就意味着存在一个双射 $f$，让我们能够“数出”集合 $A$ 的元素：第 $0$ 个元素是 $f(0)$，第 $1$ 个是 $f(1)$，……第 $n$ 个是 $f(n)$……\footnote{是的，我们从 $0$ 开始数。当然也可以从~$1$ 开始。这不会有太大区别，我们只需将~$\Nat$ 替换为~$\PosInt$ 即可。}补充以下定义可以使这一直观更加清晰：
\end{explain}

\begin{defn}
  \ollabel{defn:enumerable}
  一个集合~$A$ 是\emph{可数的}，当且仅当 $A = \emptyset$ 或者存在~$A$ 的一个枚举。我们称 $A$ 是\emph{不可数的}，当且仅当 $A$ 不是可数的。
\end{defn}

\begin{explain}
所以，一个集合是可数的，当且仅当它是空集，或者可以用枚举数出它的元素。
\end{explain}

\begin{ex}
枚举自然数集的函数就是恒等函数 $\Id{\Nat} \colon \Nat \to \Nat$，由 $\Id{\Nat}(n) = n$ 给出。枚举\emph{正}自然数集 $\Nat^+ =
\Nat \setminus \{0\}$ 的函数是 $g(n) = n + 1$，即后继函数。
\end{ex}

\begin{prob}
证明一个集合 $A$ 是可数的，当且仅当 $A = \emptyset$ 或者存在一个满射 $f\colon \Nat \to A$。证明 $A$ 是可数的，当且仅当存在一个单射 $g\colon A \to \Nat$。
\end{prob}

\begin{ex}
由
\begin{align*}
  f(n) & = 2n \text{ 和}\\
  g(n) & = 2n+1
\end{align*}
给出的函数 $f\colon \Nat \to \Nat$ 和 $g \colon \Nat \to \Nat$ 分别枚举了偶自然数集和奇自然数集。但两者都不是满射，所以都不是 $\Nat$ 的枚举。
\end{ex}

\begin{prob}
为平方数 $1$, $4$, $9$, $16$, \dots 定义一个枚举。
\end{prob}

\begin{ex}
令 $\lceil x \rceil$ 为\emph{向上取整}函数，即将 $x$ 向上取整为最近的整数。则由下式给出的函数 $f \colon \Nat \to \Int$：
\[
  f(n) = (-1)^{n} \left\lceil\tfrac{n}{2}\right\rceil
\]
如下枚举了整数集~$\Int$：
\[
\begin{array}{c c c c c c c c}
f(0) & f(1) & f(2) & f(3) & f(4) & f(5) & f(6) & \dots \\ \\
\big\lceil \tfrac{0}{2} \big\rceil & -\big\lceil \tfrac{1}{2}\big\rceil &  \big\lceil \tfrac{2}{2} \big\rceil & -\big\lceil \tfrac{3}{2} \big\rceil & \big\lceil \tfrac{4}{2} \big\rceil  & -\big\lceil \tfrac{5}{2}\big\rceil & \big\lceil \tfrac{6}{2} \big\rceil & \dots \\ \\
0 & -1 & 1 & -2 & 2 & -3 & 3& \dots
\end{array}
\]
注意 $f$ 是如何通过在正负整数之间来回“跳跃”来生成 $\Int$ 的值的。你也可以把 $f$ 看作按如下分情况定义的函数：
\[
f(n) = \begin{cases}
  \frac{n}{2} & \text{如果 $n$ 是偶数}\\
  -\frac{n+1}{2} & \text{如果 $n$ 是奇数}
  \end{cases}
\]
\end{ex}

\begin{prob}
证明如果 $A$ 和 $B$ 都是可数的，那么 $A \cup B$ 也是可数的。
\end{prob}

\begin{prob}
通过对 $n$ 进行归纳来证明：如果 $A_1$, $A_2$, \dots, $A_n$ 都是可数的，那么 $A_1 \cup \dots \cup A_n$ 也是可数的。
\end{prob}

\end{document}